\begin{table}[!htbp]\centering\tabsize
\caption{\Bi{Acquisition and selection counts by record type.}{按记录类型统计的获取与选择数量。}}\label{tab:acquisition}
\begin{tabularx}{\linewidth}{@{}p{0.24\linewidth}p{0.21\linewidth}Y@{}}\toprule
\Bi{Set}{集合} & \Bi{Count}{数量} & \Bi{Interpretation}{解释}\\\midrule
\Bi{Database query hits}{数据库检索命中} & 216 / 75 / 13 / 10 & \Bi{Scopus / WoS / IEEE / ACM; queries overlap}{Scopus / WoS / IEEE / ACM；检索间存在重复}\\
\Bi{arXiv query families}{arXiv 查询组} & 12 / 8 & \Bi{Q1--Q3: 3 + 6 + 3 hits / 8 distinct records}{Q1--Q3：3 + 6 + 3 条命中／8 条唯一记录}\\
\Bi{Focused retrieval}{定向获取} & 6 + 1 / 3 & \Bi{Six search records plus one targeted record; three sources used in qualitative comparisons}{六条检索记录及一条定向记录；三个用于定性比较的来源}\\
\Bi{Candidate inventory}{候选清单} & 155 & \Bi{Work families after linking related versions}{相关版本归组后的工作家族}\\
\Bi{Selected from inventory}{清单内采用} & 34 & \Bi{One coded source per family}{每个家族提供一个编码来源}\\
\Bi{Not selected}{未采用} & 113 + 8 = 121 & \Bi{113 not selected for claims; 8 contextual candidates not selected for core comparison}{113 项未用于主张提取；8 项背景候选未用于核心比较}\\
\Bi{Sources outside inventory}{清单外来源} & 21 & \Bi{Additional coded sources}{另行采用的编码来源}\\
\Bi{Coded subset}{编码子集} & 34 + 21 = 55 & \Bi{33 primary mechanisms and 22 auxiliary sources; 75 claims}{33 个主要机制来源及 22 个辅助来源；75 条主张}\\
\bottomrule\end{tabularx}
\end{table}
